\section{SISTEMA TERMODINÁMICO}
  Un \textit{sistema} es una región del espacio elegida para su análisis, delimitada del alrededor o \textit{entorno} por una \textit{frontera}; la frontera que separa al sistema del entorno puede ser real o imaginaria, estática o móvil \parencite{bib:cengel2015termodinamica}. Para \textcite{bib:castellan1998fisicoquimica}, un sistema termodinámico es aquella parte del universo físico cuyas propiedades se están investigando. Para el propósito de este proyecto, a partir de aquí en adelante, todo sistema termodinámico será referido simplemente como sistema. 

  Para \parencite[p.28]{bib:rodriguez2010termodinamica}, "los sistemas termodinámicos quedan determinados por sus \textit{propiedades fisicoquímicas}, los campos externos que actúan sobre ellos y las propiedades de sus paredes, tanto si esta los separan del medio exterior como si son de separación entre sus distintos subsistemas."

  Existen tres tipos de sistemas:

  \begin{enumerate}[a)]
    \item Sistemas abiertos, cuando la masa pasa a través de la frontera.
    \item Sistemas cerrados, cuando no hay transferencia de masa entre el sistema y el entorno; pero si existe transferencia d energía.
    \item Sistemas aislados, cuando la frontera evita cualquier interacción con el entorno o medio exterior.
  \end{enumerate}

  \figTiposSistema{fig:TiposSistema}

\section{SISTEMAS HOMOGÉNEOS}
  \Parencite{bib:rodriguez2010termodinamica} define un sis XXXSe denomina sistema homogéneo a aquel cuya composición química y propiedades físicas macroscópicas locales son iguales en todos sus puntos, mientras que se entiende por sistema heterogéneo al constituido por varios subsistemas homogéneos o fases.

\subsubsection{CANTIDADES Y UNIDADES}


\section{MEDIDA Y MEDICIÓN}
  Para \parencite{bib:} la palabra medida puede tener distintos significados y la palabra 'medición' define como la acción de medir.

  El resultado de la medida 
\section{SISTEMA DE CANTIDADES}
  

\section{MAGNITUDES Y UNIDADES}
  \label{sec:magnitudes_unidades}

  Según \parencite{bib:AlvarezHuayta2014}, una magnitud es aquella propiedad factible de ser cuantificada; define \textit{medición} a la técnica experimental empleada para la obtener el valor de la magnitud, y el valor obtenido a través de la medición se llama \textit{medida}. Para \parencite{bib:siu_spain2019}, al fijar el valor numérico exacto, resulta definida también la unidad. % la existencia del valor numérico está sujeto a la unidad de la medida

  \parencite[p. 2]{bib:siu_newell2019the} de \gls{nist}, adopta una perspectiva diferente; no define propiamente a la magnitud, pero si enlaza el valor que puede adoptar una medida en relación a una unidad de referencia. También define la cantidad (entendiéndose como una cantidad física) como una propiedad medible que se puede describir mediante un número y una unidad. La unidad de una cantidad es como una cantidad definida por convención que se usa como referencia para medir otras cantidades del mismo tipo. y una medida como el proceso experimental mediante el cual se obtiene el valor de una cantidad en función de una unidad.


  Por lo tanto, de acuerdo a la \parencite{bib:AlvarezHuayta2014} y 
  en la tabla~\ref{tbl:unidades_sistema_internacional}, una magnitud está compuesto 
  de los siguientes atributos: Un nombre (longitud, tiempo, masa, etc.), un 
  símbolo que representa a la magnitud, un valor numérico y una unidad de la medida. 

  Las magnitudes físicas, de acuerdo al \href{https://www.bipm.org/en/measurement-units}{SI}, se dividen en magnitudes fundamentales y derivadas. Las magnitudes fundamentales son aquellas que no se pueden expresar en términos de otras magnitudes, mientras que las magnitudes derivadas son derivaciones de las magnitudes fundamentales.
      % \newacronym{cgpm}{CGPM}{Conferencia General de Pesas y Medidas}

\section{CLASIFICACIÓN DE LAS MAGNITUDES}
  Las magnitudes se pueden clasificar en dos grupos
    \begin{enumerate}
      \item Según su origen 
        \begin{itemize}
          \item Magnitudes fundamentales
          \item Magnitudes derivadas
          \item Magnitudes suplementarias o auxiliares
        \end{itemize}
      \item Según su naturaleza 
        \begin{itemize}
          \item Magnitudes escalares
          \item Magnitudes vectoriales
          \item Magnitudes tensoriales
        \end{itemize}
    \end{enumerate}
    
  \subsection{MAGNITUDES FUNDAMENTALES}
  Son siete las magnitudes fundamentales del SI, que son la base para definir todas las demás magnitudes físicas (ver tabla~\ref{tbl:unidades_sistema_internacional}):

  \tblUnidadesSI{tbl:unidades_sistema_internacional}

  \subsection{MAGNITUDES DERIVADAS}
  Las unidades derivadas son parte del \gls{si} So aquellas que están expresadas en base a las magnitudes fundamentales; es decir, se derivan de las magnitudes fundamentales.
  
  \tblUnidadesDerivadas{tbl:unidades_derivadas_sistema_internacional}

  Las expresiones para una magnitud (o cantidad, según \gls{nist}) pueden escribirse de las formas. Los símbolos de las unidades son tratados como entidades matemáticas; Por ejemplo, la ecuación $p = 25$ kPa, o de la forma $p/$kPa = $25$; la última expresión corresponde a la forma de cociente de la cantidad y su unidad \parencite[p. 30]{bib:siu_newell2019the}. % es una cita de texto: https://www.scribbr.es/category/normas-apa/

  \fcolorbox{blue}{blue!30}{Aquí añadir una tabla y una gráfica de cómo usar las nomen} \marginpar{Aum}



  % https://www.rae.es/dpd/ordinales
  En Mayo de 2019, en la 11.\textsuperscript{a}  \gls{cgpm}, aprobó una nueva definición de las unidades del SI, que entró en vigor el 20 del mismo mes; las nuevas definiciones se basa en constantes fundamentales de la naturaleza, lo que permite una mayor precisión y estabilidad en las mediciones. 

  Las constantes fundamentales son definidas por \gls{cgpm} en la reunión del año 2019 (París, Francia), son:

  \begin{itemize}
    \item la frecuencia de transición hiperfina no perturbada en estado básico del átomo de cesio-133 $\Delta \nu_{Cs}$ es de $9 192 631 770$ Hz
    \item la velocidad de la luz en el vacío $c$ es de $299 792 458$ m/s
    \item la constante de Planck $h$ es de $6.626 070 15$ x $10^{-34}$ J s
    \item la carga elemental $e$ es de $1. 602 176 634$ x $10^{-19}$ C
    \item la constante de Boltzmann $k$ es $1.380 649$ x $10^{-23}$ J/K
    \item la constante de Avogadro $N_A$ es $6.022 140 76$ x $10^{23}$ mol\textsuperscript{-1}
    \item la eficacia luminosa de la radiación monocromática de frecuencia $540$ x $10^{12}$ Hz, Kcd, es 683 lm/W
  \end{itemize}

  % \parencite{bib:Serway2014}.
  \section{CARACTERÍSTICA DE UNA MAGNITUD} % PROPIEDADES o PROPIEDADES MAGNITUDES O ELEMENTOS
    Una magnitud está compuesta por algunas características que permiten distinguirla de forma única unas de otras:
    \begin{itemize}
      \item símbolo, es la representación gráfica de la magnitud. Ejemplos: \textbf{$m$}, $\rho$, $t$, etc.
      \item nombre, es el concepto físico de la magnitud Ejemplos: \textbf{$m$}, para al masa; $\rho$, para la densidad; $t$ para el tiempo y asi sucesivamente.
      \item valor, es valor numérico que adquiere la magnitud al momento de la medición o a través de una operación matemática.
      \item unidad, es la característica que le da sentido físico a la magnitud.
    \end{itemize}


  % Las magnitudes físicas se pueden clasificar en dos tipos:
  %   \subsubsection{Magnitudes escalares} Son aquellas que se describen completamente con un número y una unidad de medida. Por ejemplo, la temperatura, la masa y el volumen son magnitudes escalares.
  %   \subsubsection{Magnitudes vectoriales} Son aquellas que requieren tanto un número como una dirección para ser descritas. Por ejemplo, la velocidad y la fuerza son magnitudes vectoriales.


  \section{MAGNITUDES EN LA TERMODINÁMICA} % MAGNITUDES O DIMENSIONES
    Una magnitud es una propiedad física medible de un fenómeno (masa, longitud, tiempo, cantidad de sustancia, etc.) 
    Una dimensión es la naturaleza o clase de magnitud
    Para \parencite{bib:thermo:VanNess} las dimensiones fundamentales de la termodinámica (con el fin de abordar problemas reales) comienza con la identificación y delimitación de un agregado particular de materia, llamado \textit{sistema}; este sistema se define a través de algunas propiedades { \color{red} mensurables como la temperatura, la presión, el volumen cantidad de materia y el tiempo }. Estas magnitudes se relacionan entre sí a través de las leyes de la termodinámica, que describen cómo se transfieren y transforman la energía y la materia en un sistema. 

\section{PROPIEDADES DE LA MATERIA}
  Toda porción de materia posee características únicas que permiten distinguir de otras y identificarlas; esta característica se llama \textit{propiedad}. También llamado \textit{variable} o \textit{parámetro} para \parencite{bib:Himmelblau1997}, es una característica medible del material; tales como la presión, volumen o temperatura o que se pueda calcular (si no se mide directamente); además las propiedades de un sistema dependen de su condición en un momento dado y no de lo que haya sucedido en el pasado. Pueden clasificarse bajo estos criterios:

  \begin{enumerate}[a)] % a), b), c), ...
    \item Por su dependencia de la cantidad de materia: Extensivas e intensivas,
    \item Por su naturaleza: Físicas y químicas,
    \item Según su composición: Sustancias puras y mezclas (homogéneas y heterogéneas).
  \end{enumerate}

  \subsubsection{Según la dependencia de la cantidad de materia}
    Las propiedades según esta clasificación se dividen en dos tipos:
    \begin{enumerate}
      \item Las \textit{propiedades extensivas} son todas aquellas que depende de la cantidad de sustancia; es decir, su valor cambia con la cantidad de sustancia; algunos de estos son: masa, volumen, peso, capacidad calorífica, entre otros.
      \item Las \textit{propiedades intensivas}, estas propiedades son independientes de la cantidad de sustancia, son invariantes; como estos: densidad, solubilidad, volumen específico, calor específico, temperatura de fusión y otros.
    \end{enumerate}


      \subsubsection{Energía}
        \parencite{bib:castellan1998fisicoquimica} define energía como una cantidad que fluye a través de la frontera de un sistema durante un cambio de estado en virtud de una diferencia de temperatura entre el sistema y el entorno y que fluye desde una región de temperatura mayor hacia otro de menor de temperatura menor;

      
      \subsubsection{Calor}
        A partir del calor específico para una cantidad de sustancia, y un cambio de temperatura en la muestra de la sustancia, podemos decir que el calor ha sido absorbido o liberado desde o hacia el entorno. La ecuación que permite calcular la cantidad de calor involucrado en el proceso está dado por;

        \eqHeatChangeWithSpecificHeat

      \subsubsection{Capacidad calorífica}
        La capacidad calorífica (\heatCapacity[]{}) de una sustancia, es la cantidad de calor necesaria para elevar en un grado Celsius ($1 \unit{\celsius}$) la temperatura de una cantidad determinada de dicha sustancia.

      \subsubsection{Calor específico}
        El calor específico ($\specificHeat[]{}$) de una sustancia, como mencionamos previamente, es una propiedad intensiva y es la cantidad de calor que se necesita para aumentar al temperatura de un gramo ($1 \unit{\gram}$) de sustancia en $1 \unit{\celsius}$, generalmente se expresa en $\unit{\joule/(\gram \cdot \degreeCelsius)}$.    

      La expresión que relaciona estas dos propiedades es;
      
      \eqHeatCapacity

  \section{PROPIEDAD TERMODINÁMICA}
  Son características físicas medibles que permiten estudiar la energía que gana o pierde un sistema

   Las propiedades termodinámicas son aquellas que no dependen de la historia anterior de la sustancia ni los medios por el cual alcanzan un estado determinado

% \subsection{Propiedades o Variables Termodinámicas}
% El estado termodinámico queda establecido mediante las propiedades termodinámicas como la temperatura, presión y densidad.
% Las cantidades intensivas que no depende a los sucesos que se haya sometido el sistema termodinámico; estas propiedades termodinámicas se denominan {\it funciones de estado} \parencite{thermo:VanNess}

      


    % Durante el proceso de flujo estacionario 

\subsection{ESTADO}
  Según \parencite[p. 391]{bib:Himmelblau1997}, es conjunto de propiedades de los materiales en un momento dado. El estado de un sistema no depende de la "forma" o de la configuración del sistema; sino, solo de sus propiedades intensivas (temperatura, presión y composición).


\section{ESTADOS DE LA MATERIA}
Todas las sustancias pueden existir en tres estados: sólido, líquido y gas... Los gases difieren de los líquidos y sólidos en la distancia que media entre las moléculas. En un sólido, las moléculas se mantienen juntas de manera ordenada, con escasa libertad de movimiento. Las moléculas de un líquido están cerca unas de otras, sin que se mantengan en una posición rígida, por lo que pueden moverse. En un gas, las moléculas están separadas entre sí por distancias grandes en comparación con el tamaño de las moléculas mismas. \parencite[p. 9]{bib:ChangQmc2010}

\figEstadosMateria{fig:EstadosMateria}

Son posibles las transiciones entre los estados de la materia sin que este implique un cambio en su composición. Los cambios estado de materia suceden de la manera en que se puede apreciar en la figura \ref{fig:pt_phase_diagram}.

Para \textcite{bib:cengel2015termodinamica}, la velocidad de las moléculas depende de la temperatura; cuanto mayor sea la temperatura, mayor será la energía cinética media de traslación de las moléculas. En la fase líquida, las fuerzas intermoleculares son más débiles y también adopta la forma del recipiente  que los contiene. Para \parencite{bib:ReyesChumacero2012}, las principales diferencias de los distintos estados de la materia están en el arreglo dinámico de sus partículas (posición, distancia y movimiento).

% :os estados de la materia es definida la por la temperatura; en el estado gaseoso prima el movimiento de traslación y la energía cinética de las moléculas depende directamente de la temperatura del mismo; 

% La fase gaseosa, donde las moléculas están bastante apartadas y sin orden aparente, adopta la forma del recipiente que los contiene \parencite[p. 111]{bib:cengel2015termodinamica}; 



\subsection{FASE}
  De acuerdo a \parencite[p. 30]{bib:thermo:VanNess}, una fase es una región homogénea de la materia; por ejemplo: una mezcla de gases, una solución líquida y un sólido cristalino son algunos ejemplos de fases. Para \parencite[p. 461]{bib:ChangQmc2010}, una fase es una parte homogénea de un sistema, y aunque está en contacto con otras partes del sistema, está separada de esas partes por un límite bien definido. para \parencite[p. 28]{bib:rodriguez2010termodinamica}, una fase es un sistema homogéneo cuyas propiedades físicas macroscópicas locales son iguales en todo sus puntos.
    
  Pueden coexistir varias fases con la condición de que deben estar en equilibrio. A los líquidos y sólidos también se les conoce como fases condensadas o (estados condensados) \parencite[p. 430]{bib:SilberberglQmcGeneral}. 

    % Para \parencite{bib:Rodrigues2003}

\subsection{CAMBIOS DE FASE}
    También conocido como transformaciones de fase, ocurre cuando se agrega o quita energía, generalmente en forma de calor. Los cambios de fase son cambios físicos que se distinguen porque  cambia el orden molecular, y en la fase gaseosa las moléculas presentan un mayor desorden \parencite[p. 489]{bib:ChangQmc2010}. Los cambios entre los diferentes estados de agregación, según \parencite{bib:ReyesChumacero2012}, se debe a las interacciones a nivel atómico y molecular de las sustancias y se manifiestan a nivel macroscópico como fuerza de {\it cohesión}, el cual sucede entre moléculas idénticas y las fuerzas de {\it adhesión} que se dan entre moléculas diferentes.

    \figPhaseDiagram{fig:pt_phase_diagram}

    \apaFigure{fig:mi_grafico}
              {Fases de transición de estados de una sustancia}
              {resources/images/vectorial/pt_phase_diagram.pdf}
              {Elaboración propia basada en datos del XXXX .}


De la figura \ref{fig:pt_phase_diagram}, la temperatura de equilibrio para la coexistencia de la fase Sólida-Líquida a una presión dada se llama temperatura de fusión normal o temperatura de congelación normal; para la temperatura de equilibrio a una determinada presión entre la fase Líquido-Gas se llama {\it Temperatura de ebullición normal}. "Si el punto triple está a una presión superior a 1 , la sustancia no tiene una temperatura de congelación normal o una temperatura de ebullición normal, pero tiene una temperatura de sublimación normal a la que el sólido y el gas coexisten a una presión igual a 1 atmósfera", \parencite[pp. 28]{bib:Mortimer2008} %47   

% \subsection{Estructura Cristalina}
%     La fase sólida se divide en dos categorías:
%     % \begin{enumerate }
%     %   \item sólidos cristalinos 
%     %   % (las fuerzas que mantienen su estabilidad pueden ser iónicas,  covalentes, de van der Waals, enlaces de hidrógeno o una combinación de todas ellas.)
%     %   \item sólidos amorfos (ej. vidrio)
%     % \end{enumerate}

%     \subsubsection*{Tipos de cristales}
%     \begin{itemize}
%       \item cristales iónicos (formado por cationes y aniones que suelen ser de distinto tamaño)
%       \item cristales covalentes (se mantienen en una red tridimensional. Ej: diamante y grafito)
%       \item cristales metálicos (los electrones están deslocalizado en todo el cristal)
%     \end{itemize}


% antes cambio de estado: No se puede hablar de equilibrio sin antes definir el concepto de estado termodinámico
\subsection{EQUILIBRIO}
  Según \parencite[p. 29]{bib:thermo:VanNess}, el equilibrio es la ausencia de cambio, lo cual denota una condición estática; bajo esta condición no puede ocurrir ningún cambio de estado en ausencia de fuerzas que pueden promover un cambio de estado.

  Existen diferentes tipos de equilibrio, de los cuales se mencionan a continuación:
  \begin{itemize}
    \item Equilibrio mecánico,
    \item equilibrio térmico,
    \item equilibrio químico,
    \item equilibrio de fases,
    \item equilibrio termodinámico.
  \end{itemize}

\subsubsection{Equilibrio mecánico} 
  para un sistema dividido en dos regiones (región A, región B); cuando la región $\regionA$ se expande reversiblemente en una cantidad infinitesimal, la región $\regionB$ se contrae en la misma cantidad infinitesimal; el volumen total debe permanecer constante; entonces, la alta presión se expande a expensas de la región de baja presión \parencite[p.~226]{bib:castellan1998fisicoquimica}. La condición de equilibrio mecánico se alcanza cuando,

  \begin{equation*}
    \pressure[]{\regionA} = \pressure[]{\regionB}
  \end{equation*}
  
  \subsubsection{Equilibrio térmico} 
    Cuando las propiedades de dos sistemas con paredes térmicamente conductoras entran en contacto, se dice que los sistemas están en contacto térmico. las propiedades de lso sistemas se estabilizan y no cambian de valor en el tiempo, podemos decir que los sistemas están en equilibrio térmico.
    
  \subsubsection{Equilibrio químico} 
    Según \textcite{bib:Perry2000}, el equilibrio químico es un estado de reacciones reversibles donde las concentraciones o las presiones de productos y reactivos permanecen constantes; lo cual significa que, las velocidades de reacción directa e inversa son iguales.

  \subsubsection{Equilibrio de fases} 
    Es el estado en el que múltiples fases (sólido, líquido, gaseoso) coexisten sin que exista un transferencia neta de materia entre ellas. Un componente puede existir en más de una fase simultáneamente en equilibrio \textcite{bib:ball2004fisicoquímica}.

  \subsubsection{Transiciones de fase}

  \figPhaseTransition{fig:fasesTransicion}

  \subsubsection{Equilibrio termodinámico} 
    Para \parencite{bib:Reisel2021} y \parencite{bib:cengel2015termodinamica}, el equilibrio termodinámico sucede cuando un sistema alcanza el equilibrio mecánico, térmico, químico y el equilibrio de fases.
    Un sistema en equilibrio termodinámico es incapaz de realizar un cambio espontaneo de su estado; cualquier cambio de estado recibe el nombre de proceso \parencite[p. 39]{bib:cengel2015termodinamica}. El equilibrio termodinámico es el estado de balance donde todas las fuerzas del sistema termodinámico se cancelan mutuamente.


\subsection{VARIABLES DE ESTADO}
Las variables de estado definen el estado de un sistema termodinámico. Para \parencite{bib:castellan1998fisicoquimica}, el estado se define mediante las siguientes propiedades: masa, volumen, presión y temperatura. Mediante una  \textit{ecuación de estado}, puede determinarse la cuarta propiedad, el cual se conoce a partir del comportamiento del sistema.

\section{ESTADO TERMODINÁMICO}
Un estado termodinámico queda definido por las \textit{propiedades termodinámicas} \parencite[26]{bib:thermo:VanNess}. Son  propiedades termodinámicas la presión, temperatura y el volumen; estas propiedades depende unos de otros.


\subsection{FUNCIONES DE ESTADO}
Según \textcite[p.~40]{bib:sala1997termodinamica}, para especificar el \textit{estado de equilibrio} de un sistema es necesario conocer de sus variables termodinámicas; estas variables están relacionadas entre sí, no es necesario conocer el valor de todas. El número mínimo de variables de estado que define el equilibrio, se denomina grados de libertad del sistema ($\freedomDegree$).



